\documentclass[a4paper]{article}
\usepackage{amsmath,amssymb,amsthm}
\usepackage{hyperref}
\newtheorem{theorem}{Theorem}[section]
\newtheorem{definition}[theorem]{Definition}
\newtheorem{proposition}[theorem]{Proposition}
\newtheorem{lemma}[theorem]{Lemma}
\newtheorem{corollary}[theorem]{Corollary}
\theoremstyle{remark}
\newtheorem{remark}[theorem]{Remark}
\newtheorem{example}[theorem]{Example}
\usepackage[depth=2]{semtex}
\semtextrack{definition,theorem,proposition,lemma,corollary,remark,example}

\begin{document}
\section{Categories}

\subsection{Data}

A category $\mathsf{C}$ consists of the following data:
\begin{enumerate}
  \item A collection $\mathrm{Ob}(\mathsf{C})$ of objects.
  \item For each pair of objects $X, Y \in \mathrm{Ob}(\mathsf{C})$,
    a set $\mathrm{Hom}_{\mathsf{C}}(X, Y)$ of morphisms from $X$ to $Y$.
  \item For each object $X$, an identity morphism
    $\mathrm{id}_X \in \mathrm{Hom}_{\mathsf{C}}(X, X)$.
  \item For each triple of objects $X, Y, Z$, a composition map
    \begin{equation}
      \circ \colon \mathrm{Hom}(Y, Z) \times \mathrm{Hom}(X, Y)
        \to \mathrm{Hom}(X, Z).
    \end{equation}
\end{enumerate}

\subsection{Axioms}

subject to the following axioms:

\begin{definition}[Category]
A \emph{category} $\mathsf{C}$ consists of the data above,
satisfying associativity and identity laws:
\begin{equation}
  h \circ (g \circ f) = (h \circ g) \circ f
\end{equation}
and
\begin{equation*}
  f \circ \mathrm{id}_X = f = \mathrm{id}_Y \circ f.
\end{equation*}
\end{definition}

\begin{proposition}
Every category has a unique identity morphism for each object.
\end{proposition}

\begin{proof}
Suppose $e$ and $e'$ are both identities for $X$.  Then
\[
  e = e' \circ e = e'.
\]
The first equality uses $e'$ as a left identity; the second
uses $e$ as a right identity.
\end{proof}

Another paragraph after the proof.

\begin{remark}
The identity morphism is unique by the argument above.
This is a standard trick in algebra.
\end{remark}

\subsection{Examples}

\begin{example}
The category $\mathsf{Set}$ has sets as objects and functions
as morphisms.  Composition is ordinary function composition:
\begin{equation}
  (g \circ f)(x) = g(f(x)).
\end{equation}
\end{example}

\begin{example}
The category $\mathsf{Grp}$ has groups as objects and group
homomorphisms as morphisms.
\end{example}

A paragraph between examples.

\begin{example}
Any partially ordered set $(P, \leq)$ defines a category:
objects are elements of $P$, and there is a unique morphism
$x \to y$ iff $x \leq y$.  Composition and identities are forced.
\end{example}

\section{Functors}

\subsection{Definition}

A functor maps one category to another, preserving composition.

\begin{definition}[Functor]
A \emph{functor} $F \colon \mathsf{C} \to \mathsf{D}$ consists of:
\begin{enumerate}
  \item An object map $F \colon \mathrm{Ob}(\mathsf{C})
    \to \mathrm{Ob}(\mathsf{D})$.
  \item A morphism map $F_{X,Y} \colon \mathrm{Hom}_{\mathsf{C}}(X, Y)
    \to \mathrm{Hom}_{\mathsf{D}}(F(X), F(Y))$.
\end{enumerate}
satisfying:
\begin{equation}
  F(\mathrm{id}_X) = \mathrm{id}_{F(X)}
\end{equation}
and
\begin{equation}
  F(g \circ f) = F(g) \circ F(f).
\end{equation}
\end{definition}

\subsection{Properties}

\begin{lemma}
Functors preserve isomorphisms.
\end{lemma}

\begin{proof}
If $f \colon X \to Y$ is an isomorphism with inverse $f^{-1}$, then
\[
  F(f) \circ F(f^{-1}) = F(f \circ f^{-1}) = F(\mathrm{id}_Y)
    = \mathrm{id}_{F(Y)}.
\]
Similarly $F(f^{-1}) \circ F(f) = \mathrm{id}_{F(X)}$.
\end{proof}

\begin{corollary}
If $\mathsf{C}$ and $\mathsf{D}$ are isomorphic categories, then
a functor $F \colon \mathsf{C} \to \mathsf{D}$ that is an equivalence
preserves all categorical structure up to isomorphism.
\end{corollary}

A final paragraph.  Cross-reference test: see
Definition~\ref{def:functor} if we had labeled it.

\end{document}
